\section*{Domain Renewal Billing Report}

As part of a domain name registration service billing system, a company needs a report containing all customer accounts, the top-level domains (TLDs) they own, and the total annual renewal cost of domains scheduled to expire in the next calendar month.

The current month is \textbf{July 2022}, therefore only domains expiring in \textbf{August 2022} should be included.

\section*{Database Schema}

\subsection*{accounts}

\begin{center}
\begin{tabular}{|c|c|}
\hline
\textbf{column} & \textbf{type} \\
\hline
id & SMALLINT \\
\hline
username & VARCHAR(255) \\
\hline
\end{tabular}
\end{center}

\vspace{0.5cm}

\subsection*{domains}

\begin{center}
\begin{tabular}{|c|c|}
\hline
\textbf{column} & \textbf{type} \\
\hline
account\_id & SMALLINT \\
\hline
name & VARCHAR(255) \\
\hline
expiration\_date & VARCHAR(19) \\
\hline
\end{tabular}
\end{center}

\vspace{0.5cm}

\subsection*{prices}

\begin{center}
\begin{tabular}{|c|c|}
\hline
\textbf{column} & \textbf{type} \\
\hline
tld & VARCHAR(5) \\
\hline
price & VARCHAR(6) \\
\hline
\end{tabular}
\end{center}

\section*{Requirements}

Generate a report containing:

\begin{itemize}
    \item \texttt{username}
    \item One column for each TLD
\end{itemize}

Each TLD column should contain:

\begin{itemize}
    \item The total renewal cost for all domains of that TLD owned by the user
    \item Values rounded to two decimal places in currency notation
\end{itemize}

\section*{Column Ordering}

TLD columns must be ordered by:

\begin{enumerate}
    \item Descending TLD price
    \item Ascending TLD name
\end{enumerate}

\section*{Row Ordering}

Sort the result by:

\begin{itemize}
    \item \texttt{username} in ascending order
\end{itemize}

\section*{Notes}

\begin{itemize}
    \item The TLD must be extracted from the domain name.
    \item Only domains expiring in August 2022 should be included.
    \item Domain renewal totals are calculated as:
\end{itemize}

\[
\text{Total Renewal Cost}
=
(\text{Number of Domains})
\times
(\text{TLD Price})
\]

\section*{Example}

\subsection*{domains}

\begin{center}
\begin{tabular}{|c|c|c|}
\hline
\textbf{account\_id} & \textbf{name} & \textbf{expiration\_date} \\
\hline
1 & google.com & 2022-08-12 \\
\hline
1 & test.org & 2022-08-20 \\
\hline
2 & sample.net & 2022-07-18 \\
\hline
\end{tabular}
\end{center}

\vspace{0.5cm}

\subsection*{prices}

\begin{center}
\begin{tabular}{|c|c|}
\hline
\textbf{tld} & \textbf{price} \\
\hline
com & 10.00 \\
\hline
org & 8.00 \\
\hline
net & 5.00 \\
\hline
\end{tabular}
\end{center}

\section*{Expected Output}

\begin{center}
\begin{tabular}{|c|c|c|}
\hline
\textbf{username} & \textbf{col1} & \textbf{col2} \\
\hline
alice & 10.00 & 8.00 \\
\hline
\end{tabular}
\end{center}

\section*{Explanation}

\begin{itemize}
    \item \texttt{google.com} contributes \$10.00
    \item \texttt{test.org} contributes \$8.00
    \item \texttt{sample.net} is excluded because it expires in July 2022
\end{itemize}